\section{Construções em Simp}

A categoria simplicial \textbf{Simp} é restritiva devido à exigência de que os morfismos preservem a ordem linear absoluta.

\subsection{Fundamentação Teórica}
\begin{itemize}
    \item \textbf{Equalizador:} É a única construção universal básica que existe sempre em \textbf{Simp}. Consiste no subconjunto de índices onde os morfismos coincidem, o qual herda naturalmente a ordem do domínio.
    \item \textbf{Limites e Colimites (Produto, Coproduto, Pullback, Pushout):} \textbf{Não existem em geral} nesta categoria. O produto cartesiano de ordens lineares gera uma grelha (ordem parcial) que não é necessariamente uma ordem linear, falhando a propriedade universal.
\end{itemize}

\subsection{Implementação Computacional}
\begin{lstlisting}[caption={Cálculo do Equalizador Simplicial}]
% f e g sao morfismos [n] -> [m] (vetores ordenados)
idx = find(f == g); % Encontra posicoes de coincidencia
if ~isempty(idx)
    k = length(idx) - 1; % Novo objeto [k]
    eq_morphism = idx - 1; % Mapeamento para base 0
end
\end{lstlisting}